\chapter{Introduction} % Chapter 1
\label{Introduction}

% =============================================================================
% INTRODUCTION (5-8 pages)
% Rubric: Relevance/Motivation (5% of Report) + Scientific Question (10% of Report)
% =============================================================================

% -----------------------------------------------------------------------------
% SECTION 1.1: Opening Hook
% -----------------------------------------------------------------------------
\section{The Origins of Order}
% TODO: Open with the puzzle of emergent social structure
%   - When autonomous agents interact in shared environments, order emerges:
%     hierarchies, coalitions, inequality, norms
%   - Classical ABM tradition showed this with simple rules (Schelling 1971, Epstein & Axtell 1996)
%   - But now agents can REASON about their situation via LLMs
%   - Core tension: does the structure of the environment determine outcomes,
%     or does HOW agents reason, WHAT they know, and WHETHER they can commit
%     matter equally?
%
% Key references:
%   \cite{schelling1971} — micro-rules → macro-patterns (foundational)
%   \cite{epstein1996} — Sugarscape: growing artificial societies
%   \cite{axelrod1984} — cooperation from self-interested agents
%   \cite{nowak2006five} — five mechanisms for cooperation evolution

% -----------------------------------------------------------------------------
% SECTION 1.2: The Mechanism Design Framework
% -----------------------------------------------------------------------------
\section{Three Inputs to Social Order}
% TODO: Introduce the theoretical framing that structures the thesis
%   - Hurwicz (1960) mechanism design: three inputs determine outcomes in any
%     social interaction:
%     (1) How agents reason/optimize (solution concept)
%     (2) What agents know about each other (type/information space)
%     (3) How agents can communicate and commit (message space)
%   - Hobbes (1651) identified the same three factors as causes of conflict
%     in the state of nature: competition, diffidence, lack of common power
%   - This thesis manipulates all three as independent variables:
%     IV1: Reasoning depth (L0-L3 prompted reasoning)
%     IV2: Information structure (visibility, memory)
%     IV3: Communication & enforcement (no-comm → cheap talk → contracts)
%
% TODO: Brief preview of why each matters
%   - IV1: Kuusela & Roy (AAMAS 2024) showed higher reasoning → more conflict in RL.
%     Does this hold for LLMs? Is it non-monotonic?
%   - IV2: Most LLM game studies give agents full information. Real social
%     systems have local observation + memory. What changes?
%   - IV3: Hobbes's solution (Leviathan) is enforcement. Does adding credible
%     commitment shift LLM agents from conflict to cooperation?

% -----------------------------------------------------------------------------
% SECTION 1.3: Research Question
% -----------------------------------------------------------------------------
\section{Research Question}
% TODO: State main RQ:
%   "How do strategic reasoning, information structure, and enforcement mechanisms
%    each shape the emergence of social order in multi-agent LLM systems?"
%
% TODO: Operationalize:
%   - Reasoning depth: L0 (reactive) → L1 (strategic) → L2 (opponent modeling) → L3 (recursive)
%   - Information: visibility radius + persistent per-agent memory vs full info
%   - Communication: no-comm → cheap talk → binding contracts
%   - Emergent social order: inequality (Gini), cooperation rate, conflict onset,
%     community structure, phase transitions
%
% TODO: Sub-questions (mapped to IVs):
%   SQ1: How does reasoning depth interact with game structure to shape emergent order?
%        (extends Kuusela & Roy from RL to LLM)
%   SQ2: How does information structure (visibility, memory) modulate these effects?
%        (Harsanyi's incomplete information in an LLM setting)
%   SQ3: Does adding communication and enforcement shift the system from Hobbesian
%        conflict toward cooperative order?
%        (the "Leviathan hypothesis")

% -----------------------------------------------------------------------------
% SECTION 1.4: Contributions
% -----------------------------------------------------------------------------
\section{Contributions}
% TODO: Explicit contributions:
%   1. Extending the Hobbesian Trap (Kuusela & Roy) from 2-agent RL to 30-agent LLM,
%      with prompt-based reasoning manipulation on a single model (Qwen 3.5-27B)
%   2. First systematic manipulation of all three mechanism design inputs
%      (reasoning, information, communication) in LLM multi-agent systems
%   3. Persistent per-agent memory architecture as experimental variable
%      (local observation + sliding window, not god-view profiles)
%   4. Evidence for phase transitions in emergent social structure that shift
%      with reasoning depth and information structure
%   5. Reasoning traces as behavioral data (not explanation) with faithfulness
%      validation via truncation tests and Thought Anchors
%   6. Methodological: prompt-based reasoning manipulation as controlled IV
%      (vs. comparing different models, which is confounded)
%
% TODO: What this is NOT:
%   - Not claiming LLMs have "real" theory of mind
%   - Not claiming traces faithfully reflect internal computation
%   - Reasoning depth as computational intervention, not psychological claim

% -----------------------------------------------------------------------------
% SECTION 1.5: Thesis Outline
% -----------------------------------------------------------------------------
\section{Thesis Outline}
% TODO: Brief roadmap (updated chapter structure — Ch4 merged into Ch3)
%   Ch 2: Literature Review — ABMs to LLMs, mechanism design × Hobbes,
%          three IVs (reasoning, information, communication), faithfulness,
%          phase transitions
%   Ch 3: Methods — game design, agent architecture + model selection,
%          independent variables, experimental design, metrics,
%          computational infrastructure (Snellius, vLLM, reproducibility)
%   Ch 4: Results — OAT parameter characterization, reasoning depth effects,
%          information/memory effects, communication effects, phase transitions
%   Ch 5: Validation — CoT faithfulness tests, K-level reasoning verification,
%          trace analysis methodology
%   Ch 6: Discussion — interpretation, Hobbesian Trap extension, limitations
%   Ch 7: Conclusion — answers to RQs, contributions, future work
